\section*{Configuration-Type Pigeonhole: corrected scope}

\noindent\textbf{Status:} The elementary statement below is proved.
A cycle-rank-to-configuration encoding and any consequent splitting of
\emph{different} rooted vertex types remain \textbf{OPEN}.

Fix $k,\Delta,R$ and a finite collection $\mathcal T(k,\Delta,R)$ of
well-defined local configuration types. Finiteness holds when these types
are evaluated on rooted radius-$R$ balls of graphs with bounded degree
and a fixed finite relational vocabulary. No claim about arbitrary
unbounded-radius global FO$^k$ types is intended.

\begin{lemma}[Finite-type pigeonhole]
Let $\mathcal C$ be a finite set of \emph{distinct} configurations with
a specified map $\tau:\mathcal C\to\mathcal T$. If
$|\mathcal C|>|\mathcal T|$, then there exist distinct
$c,c'\in\mathcal C$ with $\tau(c)=\tau(c')$.
\end{lemma}
\begin{proof}
If $\tau$ were injective, the cardinality inequality would be reversed.
\end{proof}

\noindent\textbf{Conditional cycle-basis application.}
If $\dim_{\mathbb F_2}Z_1(B_{2R}(v))>|\mathcal T|$ \emph{and}
there is a proved injection from a basis of that cycle space into
$\mathcal C$ with a specified type map $\tau$, then the lemma gives two
\emph{distinct configurations of equal type}. The injection cannot be
inferred from the dimension inequality alone.

\noindent\textbf{Invalid former inference:} Equal types for two
configurations do not imply \emph{unequal} rooted vertex types at some
pair of vertices. Vertex-transitive graphs, including the Petersen graph,
can have nonzero local cycle rank with only one rooted vertex type.
Therefore the previous ``Consequently'' clause asserting a vertex-type
split is withdrawn; it cannot support a cycle--orbit splitting theorem.
